\section{Diskussion und Grenzen}

\subsection{Technische Machbarkeit und Robustheit}

Die Hauptfrage kann für den kontrollierten Alurohrdatensatz und den
nachgewiesenen Matrix-/Smoke-/Replaypfad grundsätzlich positiv beantwortet
werden. Die Verarbeitungskette kann aus einer Bildsequenz Masken und Kameras
erzeugen, objektbezogene Gaussians trainieren, ein Mesh ableiten und daraus
eine lokale Centerline samt B-Spline exportieren. Der wissenschaftlich
relevante Befund ist dabei nicht nur ein einzelner erfolgreicher Lauf. Die im
Matrix-/Replaypfad wirksamen Quality-Gates für Masken, Kameramodelle, ideale
Bilddomäne, Eval-Split und Ergebnisarchive verhindern mehrere zuvor
beobachtete stille oder verspätete Fehler; Inline-Vollauf, Matrixrunner und
Replays mit frischem COLMAP nutzen dieselbe Gate-/Warp-Kette, während
Replays ab \texttt{sts} einen bereits passenden idealen Maskensatz voraussetzen.

Robustheit bedeutet in dieser Arbeit nicht Fehlerfreiheit unter jeder
denkbaren Eingabe. Sie bedeutet, dass bekannte Voraussetzungen geprüft,
unvollständige Routen eindeutig markiert, teure Stufen an kontrollierten
Grenzen wiederholt und Ergebnisse reproduzierbar archiviert werden. Die
Untersuchung nur eines Testvideos begrenzt weiterhin die externe Validität.

\subsection{Masken, Domänen und Geometrie}

Die zentrale Domänentrennung löst einen methodischen Widerspruch: SfM profitiert
von unmaskierten Hintergrundmerkmalen, während STS und objektmaskierte
Metriken eine semantisch begrenzte ideale Bilddomäne benötigen. Die gemeinsame
Transformation von Bildern, Kameras und Masken stellt diese Konsistenz her.

Masken sind dennoch keine geometrische Referenz. Sie gewichten oder selektieren
Bildpixel und Objektzuordnungen. Position, Skalierung und Orientierung der
Gaussians, Kamerafehler, Surface-Sampling und Poisson-Rekonstruktion können die
3D-Geometrie weiterhin verändern. Deshalb ist die zunächst plausible Aussage,
eine niedrige Auflösung genüge bei korrekten Masken und registriertem COLMAP,
nur als technische Hypothese zulässig. Für einen erfolgreichen Lauf müssen alle
Datenverträge gelten; für geometrische Genauigkeit wäre zusätzlich eine
unabhängige Referenz notwendig.

\subsection{Kameramodell, FPS und Auflösung}

Die Matrix zeigt, dass Kameramodell, zeitliche Samplingdichte und Auflösung
nicht isoliert interpretiert werden dürfen. Die 2-FPS-Profile enthalten weniger
zeitlich benachbarte Ansichten; bei dünnen Objekten kann dies die SAM-
Propagation und die verfügbare Überdeckung erschweren. Die Maskensegmentierung
selbst war im Testdatensatz stabil: Das Rohr wurde über alle Auflösungen und
FPS-Profile hinweg in jedem Frame erkannt; instabil waren allenfalls die
numerischen Metrikwerte über die Konfigurationen hinweg (dort besteht keine
monotone Beziehung zu FPS oder Auflösung, siehe
Abschnitt~\ref{sec:visueller-aufloesungsvergleich}).

Die nativen Bildmetriken werden in unterschiedlichen Pixelräumen berechnet.
Downsampling wirkt dabei wie ein Tiefpassfilter: Feine Hochfrequenzfehler wie
Metallreflexionen oder feine Kantenverschiebungen werden beim Herunterrechnen
weggeglättet, sodass Low-Auflösung künstlich höhere Werte erreichen kann. Für
BIM und Vermessung ist 720p trotzdem überlegen: Die Objektdetailgröße pro
Pixel ist in 720p deutlich feiner als in low, sodass kleine Strukturmerkmale
des Rohrs erhalten bleiben, die im low-Bild unterschritten werden. Eine echte
GSD-Aussage wäre erst mit einer Maßstabsreferenz zulässig. Die Konsequenz ist
zweifach: Auflösungen dürfen nicht in einer globalen Rangliste verglichen
werden; belastbar sind nur gepaarte Vergleiche von STS gegen SuGaR innerhalb
derselben Auflösung. Ein strengerer
Auflösungsvergleich müsste alle Renderings auf eine gemeinsame Zielauflösung
und identische Eval-Views bringen. Die derzeitige Matrix ist deshalb ein
Engineering-Screening.

Wissenschaftlich belastbar sind deshalb nur gepaarte Vergleiche innerhalb
derselben Auflösungs-, Kameramodell- und FPS-Konfiguration. Die Boxplot-
Auswertung ergänzt die Mittelwerte um die Streuung der einzelnen Testansichten;
sie ersetzt aber nicht die geometrische Referenzmessung.

Die Bildtafeln des Pipeline-Durchlaufs in
Anhang~\ref{app:durchlauf} präzisieren diese Einschränkung: Trotz des gleichen
SIFT-Merkmalslimits sind in den Punktwolken kaum Auflösungsunterschiede
sichtbar; die sichtbaren Floater-, Rand- und Centerline-Unterschiede zeigen
jedoch, warum Bildmetrik und Geometrie getrennt bewertet werden müssen.

Der direkte PINHOLE-Lauf ist eine Ablation und kein Nachweis einer gelungenen
Entzerrung: Er erzwingt auf den Rohbildern die Annahme einer idealen Kamera,
während der Produktionsweg mit \texttt{OPENCV} zunächst Verzeichnung schätzt
und daraus eine separate ideale Szene erzeugt
(Abschnitt~\ref{sec:produktionskonfiguration}). Die von COLMAP ausgegebenen
Verzeichnungskoeffizienten sind geschätzte Parameter, keine gemessenen
Objektivdaten; der Reprojektionsfehler belegt interne Konsistenz, nicht die
physikalische Richtigkeit der Entzerrung. Für einen Vermessungsnachweis wäre
eine externe Kamerakalibrierung erforderlich; diese liegt nicht vor.

\subsection{Original-GS und SuGaR-Coarse}

Die Vierfeld-Ablation zeigt, dass nicht nur die verwendete Tiefenroute, sondern
auch die zusätzliche SuGaR-Coarse-Optimierung die resultierende Oberfläche
verändert. Route A bewahrt den STS-Zustand am direktesten und ist für das Ziel
einer stabilen Centerline deshalb der plausibelste Produktionsstandard. Mehr
Vertices oder sichtbar feinere Strukturen sind nicht automatisch besser, wenn
sie zugleich Spitzen oder veränderte Randstützen erzeugen.

Die Entscheidung ist bewusst vorläufig. Die Meshdistanzen vergleichen nur
Varianten untereinander. Sie geben keinen Abstand zum realen Alurohr an. Die
zwischenzeitlich abgeschlossene zwölfteilige SuGaR-Coarse-Folgematrix zeigt,
dass der reparierte Coarse-Renderpfad über beide Kameramodelle, drei
Auflösungen und zwei FPS-Profile technisch durchläuft. Sie liefert damit eine
deutlich breitere technische Evidenz als die frühere Einzelprüfung. Da in keinem
der zwölf Läufe ein `refined.ply` exportiert wurde, ist die Aussage weiterhin
auf SuGaR-Coarse beschränkt; eine SuGaR-Refined-Auswertung bleibt offen.

Der gepaarte Delta-Vergleich der zwölf Follow-up-Konfigurationen ist deshalb
als Rendervergleich zu lesen: SuGaR-Coarse wird jeweils mit dem STS-Splat
derselben Konfiguration verglichen. Über alle zwölf gepaarten Konfigurationen
liegt die SuGaR-Coarse-Bildmetrik durchgängig unter der STS-Baseline. Diese
konsistente Richtung ist eine zulässige Aussage der Delta-Grafik – SuGaR
schneidet in den Ansichtsmetriken durchgängig schlechter ab. Sie bleibt
allerdings ein Rendervergleich und kein Meshgenauigkeitsnachweis; die
Unterschiede können neben der zusätzlichen SuGaR-Optimierung auch aus dem
unterschiedlichen Renderpfad entstehen.

\subsection{Autopilot und Batchfähigkeit}

Der Matrixrunner belegt die serielle Batchverarbeitung deklarierter
Konfigurationen. Der Autopilot füllt interaktive Entscheidungen mit
Standardwerten und kann den manuellen GCP-Breakpoint überspringen. Beide
Funktionen sind verwandt, aber nicht identisch.

Die nutzerseitige Autonomie ist durch archivierte Autopilot-Volläufe belegt
(Abschnitt~\ref{sec:versuchsaufbau}, Erfolgskriterium~6). Die
Containerisierung und die nichtinteraktive Matrix bereiten einen Serverbetrieb
vor; ein Deployment, Mehrbenutzerbetrieb und Betriebsmonitoring wurden nicht
evaluiert. Eine formale Usability-Messung wurde ebenfalls nicht durchgeführt;
im Sinne der Nutzerfreundlichkeit wurden durch Tests gewonnene Standardwerte
jedoch so gesetzt, dass dem Nutzer wiederkehrende Entscheidungen abgenommen
werden.

\subsection{Laufzeit und Effizienz}

Der Zeit-/Qualitätsquotient ist ebenfalls sekundär. Er kombiniert normalisierte
PSNR-, SSIM- und LPIPS-Werte mit gleichen Gewichten; bei LPIPS kehrt sich das
Ranking um, da hier niedrigere Werte besser sind – das ist die Natur der
Metrik, keine Inversion der Daten. Er verwendet
nur die archivierten STS-bis-Postprocess-Zeiten. SAM~3.1- und COLMAP-Zeit sowie
Container-Overheads der Rendering- und Metrikschritte sind nicht vollständig
enthalten. Vollständige End-to-End-Zeiten liegen zwischenzeitlich aus dem
Verifikationsbatch \path{matrix_e2e_verifikation_260826} vor
(Abschnitt~\ref{sec:laufzeiten-e2e}); sie ergänzen die Ressourcenplanung,
ändern aber nichts an der Segmentbasis des Quotienten. Für die
Bachelorarbeit müssen Rohmetriken und geometrische
Referenzmetriken höher gewichtet werden.

Die Effizienzanalyse ist trotzdem praktisch relevant: Ein automatischer
Screeninglauf kann ungeeignete Kombinationen aussortieren, bevor eine
hochaufgelöste Produktionsrechnung erfolgt. Die Zeit-/Qualitätsauswertung
dient damit als Screening-Hilfe, nicht als primäre wissenschaftliche
Rangliste.

\subsection{Centerline, Georeferenzierung und externer Geltungsbereich}

Die größte technische Restaufgabe ist die echte Georeferenzierung. Ein
Translation-Fallback prüft zwar die Ausführung, sagt aber nichts über Lage-
oder Höhengenauigkeit aus. Ebenso sind die Testdaten keine reale Rohr- oder
Kabeltrasse und enthalten keine unabhängige GNSS-Referenzkurve.

Auch die lokale Centerline ist bisher nur als erzeugtes Artefakt validiert,
nicht gegen eine geometrisch definierte Referenz wie die wahre Scheitelachse
des Rohrs. Centerline und Scheitelachse fallen nur bei einer idealen
Rekonstruktion zusammen; ihr Vergleich setzt eine unabhängige Vermessung
voraus. Der \texttt{single}-Modus ist für eine einzelne Trasse robust,
verwirft aber bewusst Netzwerktopologie. Das ist für den Untersuchungsgegenstand
gewollt: Die Bachelorarbeit untersucht bewusst lineare Objekte ohne
Ringschluss oder Abzweigungen; eine Übertragung auf verzweigte Netze würde
neue Referenzdaten und eigene Graphmetriken erfordern und ist nicht Teil des
Scopes.

Für die Bachelorarbeit folgen daraus vier priorisierte Schritte: reale
Aufnahmebedingungen, echte GCP-/GNSS-Daten, validierte relative
4x4-Transformation sowie RMSE- und Hausdorff-Distanz einer extrahierten
Centerline gegen eine unabhängige Referenzkurve. Erst damit kann die
ingenieurtechnische Toleranz geprüft werden.

\subsection{Beantwortung der Forschungsfragen}

\begin{enumerate}
  \item Die konsistente Übergabe erfordert ein erhaltenes Originalmodell, eine
	  getrennte ideale Bilddomäne, denselben Warp für Masken, einen
	  nichtleeren gematchten Eval-Split und explizite Statusprüfungen aller
	  Ergebnisstufen.
  \item Kameramodell, FPS und Auflösung verändern die Ansichtsmetriken in
	  unterschiedlichen Richtungen: Kein Faktor führt über alle Konfigurationen
	  hinweg zu durchgängig besseren Werten (keine monotone Beziehung zwischen
	  Metrikwert und Faktorstufe). Die Maskensegmentierung selbst war dagegen
	  stabil – im Testvideo wurde das Rohr in jedem Frame erkannt. Die Matrix
	  erlaubt damit ein Screening, aber keinen isolierten Kausal- oder
	  Geometrienachweis.
  \item Route A bewahrt die Original-STS-Gaussians direkter und zeigt im
	  relativen Vergleich weniger unerwünschte Oberflächenveränderung. Ihre
	  absolute geometrische Richtigkeit bleibt offen.
  \item Der normale Inline-Vollauf und der Matrixrunner sind als
	  reproduzierbare Ausführungswege durch archivierte Läufe nachgewiesen
	  (Abschnitt~\ref{sec:versuchsaufbau}); der Autopilot ist über sechs
	  archivierte Volläufe belegt.
\end{enumerate}
